%Resumen del informe en no m\'as de 10 l\'ineas, donde se sintetice el problema que se trata y sirva para que un lector no involucrado comprenda el objetivo del documento.
El \textit{Truck Drone Parcel Delivery Problem} (TDPDP) aborda la distribuci\'on coordinada de paquetes mediante un cami\'on y un dron que opera en paralelo para minimizar el tiempo total de entrega. En este trabajo se analiza el problema y se presenta un algoritmo evolutivo h\'ibrido bajo el enfoque \textit{Truck-First, Drone-Second}. La propuesta optimiza recorridos base (TSP) mediante operadores gen\'eticos y asigna vuelos factibles del dron utilizando programaci\'on din\'amica (\textit{SPLIT}). La evaluaci\'on experimental sobre instancias de 10 a 200 clientes demuestra que la inicializaci\'on con GRASP mejora notablemente la calidad frente a un muestreo aleatorio. Asimismo, se evidencia que un equilibrio entre cruzamiento y mutaci\'on favorece la convergencia, y que priorizar m\'as generaciones sobre la poblaci\'on ($n < g$) optimiza la calidad y el tiempo de c\'omputo.
